\phantomsection
\addcontentsline{toc}{subsection}{Inciso b)}
\subsection*{Inciso b)}

% Seccion con nombres de clase largos, que no se parten con guion.
\sloppy

\textit{Implementa con 6 componentes y 3 actividades: VerPelícula, EscucharMúsica,
JugarVideojuego.}

\vspace{0.4cm}

\subsubsection*{Qué se implementó}

La implementación completa, compilable y ejecutable, está en
\texttt{docs/CP4/codigo/ejercicio7/} (C\# 12 sobre .NET 8, sin dependencias externas).

\begin{itemize}[leftmargin=1.2cm]
    \item \textbf{El subsistema, seis componentes:} \texttt{Amplificador},
          \texttt{ReproductorDvd}, \texttt{ConsolaJuegos}, \texttt{Proyector},
          \texttt{Pantalla} y \texttt{Luces}, sobre \texttt{ComponenteBase}, que guarda el
          nombre, el estado y la traza de lo que hace cada uno.
    \item \textbf{La fachada:} \texttt{HomeTheaterFacade}, que recibe los seis al
          construirse y publica \texttt{VerPelicula}, \texttt{TerminarPelicula},
          \texttt{EscucharMusica} y \texttt{JugarVideojuego}.
    \item \textbf{El cliente:} \texttt{Demo}, con las cinco comprobaciones cuya salida se
          reproduce más abajo.
\end{itemize}

El enunciado nombra cinco componentes y pide seis. El añadido es \texttt{ConsolaJuegos}, y no
es de relleno: sin él, \texttt{JugarVideojuego} sería una película con otro nombre.

\subsubsection*{Las tres actividades sobre los seis componentes}

Lo que hace interesante al ejercicio es que las tres actividades no son la misma secuencia con
otros parámetros: usan \emph{subconjuntos distintos} del subsistema, en órdenes distintos.

\begin{table}[h!]
    \centering
    \small
    \begin{tabular}{p{2.3cm}p{3.9cm}p{3.5cm}p{3.9cm}}
        \toprule
        \textbf{Componente} & \textbf{VerPelícula} & \textbf{EscucharMúsica}
        & \textbf{JugarVideojuego} \\
        \midrule
        Luces        & atenúa al 10\% & atenúa al 40\% & atenúa al 60\% \\
        \addlinespace
        Pantalla     & baja & --- & baja \\
        \addlinespace
        Proyector    & enciende, entrada dvd, formato 16:9 & ---
                     & enciende, entrada consola, formato 16:9 \\
        \addlinespace
        Amplificador & enciende, entrada dvd, envolvente, volumen 7
                     & enciende, entrada dvd, estéreo, volumen 5
                     & enciende, entrada consola, volumen 9 \\
        \addlinespace
        DVD          & enciende, reproduce & enciende, reproduce & --- \\
        \addlinespace
        Consola      & --- & --- & enciende, carga el juego \\
        \bottomrule
    \end{tabular}
\end{table}

Las dos celdas vacías de \texttt{EscucharMúsica} son las que más dicen: encender el proyector
para oír música sería un error que nadie comete dos veces, y es precisamente el tipo de
conocimiento que la fachada guarda.

\subsubsection*{La fachada en el código}

{\small
\begin{verbatim}
public sealed class HomeTheaterFacade
{
    private readonly Amplificador _amplificador;
    private readonly ReproductorDvd _dvd;
    // ... los otros cuatro componentes

    public void VerPelicula(string titulo)
    {
        _luces.Atenuar(10);
        _pantalla.Bajar();
        _proyector.Encender();
        _proyector.Entrada("dvd");
        _proyector.Formato("16:9");
        _amplificador.Encender();
        _amplificador.Entrada("dvd");
        _amplificador.ModoEnvolvente(true);
        _amplificador.PonerVolumen(7);
        _dvd.Encender();
        _dvd.Reproducir(titulo);
    }

    public void TerminarPelicula()   // el orden inverso tambien
    {                                // es conocimiento del subsistema
        _dvd.Detener();
        _dvd.Apagar();
        _amplificador.Apagar();
        _proyector.Apagar();
        _pantalla.Subir();
        _luces.Encender();
    }
}
\end{verbatim}
}

Obsérvese que el cuerpo no contiene ni una decisión: ni un \texttt{if}, ni un cálculo. La
fachada no añade comportamiento, añade \emph{orden}. Si alguna vez apareciera lógica de
negocio aquí dentro, sería señal de que la clase ha dejado de ser una fachada.

\subsubsection*{Qué cambia para el cliente}

\begin{table}[h!]
    \centering
    \small
    \begin{tabular}{p{4.6cm}p{4.6cm}p{4.6cm}}
        \toprule
        & \textbf{Sin fachada} & \textbf{Con fachada} \\
        \midrule
        Llamadas para ver una película & 11 & 1 \\
        \addlinespace
        Clases que el cliente conoce & 6 & 1 \\
        \addlinespace
        Quién sabe el orden & cada cliente, por su cuenta & un solo archivo \\
        \addlinespace
        Coste de cambiar el amplificador & todos los clientes & un archivo \\
        \bottomrule
    \end{tabular}
\end{table}

\subsubsection*{Evidencia de ejecución}

Salida literal de \texttt{Demo}. Las partes 1 y 2 producen deliberadamente la misma traza: lo
que cambia no es lo que ocurre, sino quién tuvo que saberlo.

{\footnotesize
\begin{verbatim}
1. El cliente SIN fachada: once llamadas y el orden es cosa suya
----------------------------------------------------------------
  codigo del cliente: once lineas, seis componentes, un orden
  que nadie ha escrito en ningun sitio.
    luces         atenuadas al 10%
    pantalla      bajada
    proyector     encendido
    proyector     entrada -> dvd
    proyector     formato 16:9
    amplificador  encendido
    amplificador  entrada -> dvd
    amplificador  modo envolvente
    amplificador  volumen 7
    dvd           encendido
    dvd           reproduciendo "Blade Runner"
    dvd           detenido
    dvd           apagado
    amplificador  apagado
    proyector     apagado
    pantalla      subida
    luces         al 100%

2. El mismo resultado CON fachada: una llamada
----------------------------------------------
  codigo del cliente: cine.VerPelicula("Blade Runner");
    luces         atenuadas al 10%
    pantalla      bajada
    proyector     encendido
    proyector     entrada -> dvd
    proyector     formato 16:9
    amplificador  encendido
    amplificador  entrada -> dvd
    amplificador  modo envolvente
    amplificador  volumen 7
    dvd           encendido
    dvd           reproduciendo "Blade Runner"

3. Terminar: el orden inverso tambien es conocimiento del subsistema
--------------------------------------------------------------------
  codigo del cliente: cine.TerminarPelicula();
    dvd           detenido
    dvd           apagado
    amplificador  apagado
    proyector     apagado
    pantalla      subida
    luces         al 100%

4. Las otras dos actividades sobre el mismo subsistema
------------------------------------------------------
  cine.EscucharMusica("Kind of Blue");
    luces         atenuadas al 40%
    amplificador  encendido
    amplificador  entrada -> dvd
    amplificador  modo estereo
    amplificador  volumen 5
    dvd           encendido
    dvd           reproduciendo "Kind of Blue"

  cine.JugarVideojuego("Hollow Knight");
    luces         atenuadas al 60%
    pantalla      bajada
    proyector     encendido
    proyector     entrada -> consola
    proyector     formato 16:9
    amplificador  encendido
    amplificador  entrada -> consola
    amplificador  volumen 9
    consola       encendida
    consola       cargando "Hollow Knight"

5. La fachada no cierra el subsistema
-------------------------------------
  Un cliente con una necesidad que la fachada no cubre sigue
  pudiendo hablar con el componente directamente:
    amplificador  volumen 3

  Estado final de los seis componentes:
    amplificador  entrada:consola
    dvd           reproduciendo
    consola       jugando
    proyector     entrada:consola
    pantalla      bajada
    luces         atenuadas
\end{verbatim}
}

\subsubsection*{Tres observaciones que salieron del experimento}

\textbf{El apagado es tan subsistema como el encendido.} \texttt{TerminarPelicula} no es la
inversa automática de \texttt{VerPelicula}: apaga en orden contrario, no devuelve el volumen
y sube la pantalla al final. Una fachada que solo supiera empezar actividades estaría a medias.

\textbf{La fachada no cierra el subsistema.} En la parte 5, el cliente baja un nivel y le pone
volumen 3 al amplificador directamente. Compila y funciona, y debe hacerlo: el catálogo lo
dice expresamente. Quien quiera impedirlo necesita otro patrón, y probablemente otro problema.

\textbf{El estado queda repartido.} Al terminar, cada componente conserva el suyo
(\texttt{entrada:consola}, \texttt{jugando}, \texttt{bajada}\ldots) y la fachada no lo
centraliza. Es coherente con el patrón, pero conviene saberlo: la fachada no es la dueña del
estado del subsistema, y preguntarle a ella en qué estado está la sala no funcionaría.

\subsubsection*{Cómo compilar y ejecutar}

{\small
\begin{verbatim}
cd docs/CP4/codigo/ejercicio7
dotnet run -- demo
\end{verbatim}
}

Compila sin un solo aviso con \texttt{TreatWarningsAsErrors} y nivel de aviso 9999 sobre
.NET 8.

\fussy
